% Fragment partage : inclus par qr-code.tex (dossier autonome)
% et par dossier-conception-alanya.tex (dossier client assemble).
% Ne pas y mettre de preambule ni de \part : c'est l'appelant qui les pose.

\chapter{Le besoin}

\section{Ce qui existe déjà}

Alanya n'a aujourd'hui \textbf{qu'un seul chemin} pour ajouter un contact préféré~: la
recherche manuelle par nom, pseudo ou Alanya ID, depuis la feuille d'ajout du profil.
Aucun lien d'invitation, aucun code, aucun geste plus rapide.

\begin{constat}[Vérifié dans le code]
\code{lib/widgets/add\_contact\_sheet.dart} interroge d'abord le cache local
(\code{filterUsersBySearch}), puis \code{GET~/users/search} après 400~ms sans résultat.
L'ajout appelle \code{POST~/contacts/:userId} --- la même route que celle que ce volet
réutilisera pour l'ajout par QR.
\end{constat}

Côté connexion, Alanya n'a \textbf{qu'un seul appareil implicite}~: se connecter sur un
second téléphone fonctionne déjà techniquement (rien ne l'empêche), mais rien ne le montre
non plus. Il n'existe aucun écran « mes appareils », aucun moyen de savoir où son compte est
ouvert, ni de fermer une session à distance.

\begin{constat}[Vérifié dans le code]
\code{lib/screens/profile/account\_security\_screen.dart} ne contient que deux
sections~: \emph{Email} et \emph{Changer le mot de passe}. Rien sur les appareils.
\end{constat}

\section{Ce qu'on ajoute}

Deux fonctionnalités, toutes deux fondées sur un QR code~:

\begin{enumerate}[leftmargin=1.6em]
  \item \textbf{Un code personnel éphémère.} Chaque utilisateur génère, en ouvrant
        « Mon code », un QR à son identité --- logo Alanya au centre, traits arrondis,
        une vraie carte de visite plutôt qu'un carré noir et blanc. Le scanner d'un autre
        utilisateur reconnaît la personne et l'ajoute \emph{instantanément} en contact
        préféré. Le code vit \textbf{dix minutes} et s'éteint au premier scan réussi~;
        l'écran en régénère un neuf tout seul, et le propriétaire est prévenu --- et
        invité à ajouter l'autre \textbf{en retour} --- chaque fois que son code est
        utilisé.
  \item \textbf{La connexion d'un second appareil.} Un nouvel appareil affiche un code ; un
        appareil déjà connecté le scanne pour l'autoriser. L'utilisateur peut ensuite voir,
        depuis son compte, la liste de tous les appareils connectés --- et en déconnecter un
        à distance.
\end{enumerate}

\begin{note}[Une seule brique, deux usages]
Les deux fonctionnalités ne partagent pas qu'un mot dans leur nom. Elles partagent un
\textbf{composant scanner}, un \textbf{générateur de QR stylé}, et une manière commune de
résoudre ce qu'un code scanné signifie. C'est pour cette raison qu'elles sont conçues dans
un seul volet --- le même raisonnement qui a réuni compte officiel et diffusion dans
l'étape~3.
\end{note}

\section{Les règles produit}

\begin{center}
\begin{tabularx}{\linewidth}{@{}p{4.6cm}X@{}}
\toprule
\thd{Qui peut scanner} & Un utilisateur \textbf{déjà connecté}, dans les deux cas. Scanner
n'est jamais un geste public ; c'est toujours un compte qui en reconnaît un autre. \\
\thd{Vie du code personnel} & \textbf{Dix minutes, une seule personne.} Généré à la
demande à l'ouverture de « Mon code », consommé au premier scan réussi. L'écran affiche
un compte à rebours et régénère automatiquement --- à l'expiration comme après
consommation. \\
\thd{Ajout de contact} & \textbf{Instantané}, sans écran de confirmation. Le filet de
sécurité est une carte de résultat annulable, pas une étape supplémentaire. \\
\thd{Retour au propriétaire} & Prévenu à \textbf{chaque scan} de son code, où qu'il soit
dans l'application, et invité à ajouter l'autre en retour --- notification si l'app est
en arrière-plan, simple information si le lien existait déjà. \\
\thd{Contacts ajoutés par QR} & Reconnaissables partout~: pastille sur l'avatar, mention
datée sur la fiche, filtre « Par QR » dans les contacts préférés \emph{et} dans la liste
des discussions. \\
\thd{Note de contexte} & Optionnelle, saisie juste après le scan --- « rencontré au salon
de Douala » --- ou au moment d'ajouter en retour, et relue sur la fiche du contact. Elle
appartient à la relation~: l'autre ne la voit jamais. \\
\thd{Codes permanents} & \textbf{Réservés aux futurs comptes business} (volet~1). Aucun
compte personnel n'a de code qui dure. \\
\thd{Connexion d'un appareil} & \textbf{Jamais instantanée.} L'appareil qui scanne voit
toujours l'identité de l'appareil qu'il autorise, et doit confirmer explicitement. \\
\thd{Durée du QR de connexion} & Environ \textbf{90~secondes}, puis régénération
automatique --- comme un code à usage unique, pas comme une carte de visite. \\
\thd{Emplacements} & Profil (mon code), feuille d'ajout de contact (scanner), Compte et
sécurité (appareils connectés), écran de connexion (nouvel appareil). \\
\bottomrule
\end{tabularx}
\end{center}

\begin{note}[Pourquoi un code qui meurt]
Un code permanent est une affiche~; un code de dix minutes est une poignée de main. Le
code personnel sert à des rencontres --- on le montre, on l'envoie, la personne en face
scanne, c'est fini. L'usage unique garantit que le code n'ajoute que la personne à qui on
vient de le tendre~: une capture d'écran qui circule ensuite ne mène plus nulle part. Et
comme l'écran régénère aussitôt, montrer son code à trois personnes d'affilée reste
trois gestes de \emph{leur} côté, zéro du sien.
\end{note}

\begin{constat}[Vérifié dans le code --- pourquoi le permanent est verrouillé]
L'infrastructure du code permanent (\code{users.qr\_public\_id}) existe et reste en
place, mais son garde refuse aujourd'hui tout le monde --- et pour une raison précise~:
la colonne \code{users.type\_compte}, qu'on prendrait volontiers pour un type de compte
métier, est en réalité un \textbf{rôle d'administration} (0~utilisateur, 1~admin,
2~super-admin, cf. la route \code{PUT~/users/:id/role} de \code{src/routes/admin.js},
côté Alanya-Backend). Rien dans la base ne sait encore dire « ce compte est un
commerce ». Le volet~1 (socle de compte) apportera \code{account\_type}~; le verrou des
codes permanents s'y branchera alors.
\end{constat}

\chapter{Mon code personnel}

\section{Vue d'ensemble}

Un seul écran, deux volets accessibles d'un geste~: \emph{Mon code} et \emph{Scanner} ---
comme on bascule entre appareil photo avant et arrière.

\begin{center}
  \imageOuCadre{maquettes/qr-identite.png}{\linewidth}%
    {Mon code (médaillon, logo au centre, partage et régénération) et l'écran Scanner
     (reconnaissance instantanée, succès doux et annulable)\\
     \texttt{maquettes/qr-identite.png}}
\end{center}

\begin{note}[Maquette interactive]
\href{https://claude.ai/code/artifact/69e9cbec-ff13-48fc-903f-673648874b2a}%
{\texttt{claude.ai/code/artifact/69e9cbec}}\\
Source versionnée~: \code{docs/architecture/maquettes/qr-identite.html}
\end{note}

\section{Écran --- Mon code}

Une carte blanche~: médaillon d'avatar débordant le haut, nom, Alanya ID, puis le QR
code lui-même --- fond clair, points arrondis à la couleur de marque, logo Alanya au
centre. Sous le code, un \textbf{compte à rebours} (« Expire dans 9:41 ») et la mention
« Valable 10~minutes et pour une seule personne ». En dessous, deux actions~:
\textbf{Partager} et \textbf{Nouveau code}.

À zéro --- comme après un scan réussi --- l'écran régénère un code neuf de lui-même~:
l'utilisateur n'a jamais un code mort sous les yeux, exactement le comportement déjà
retenu pour le QR de connexion (chapitre~\ref{ch:qr-naif-connexion}). « Nouveau code »
reste disponible pour qui veut couper court sans attendre~; aucun dialogue de
confirmation, renouveler un code qui expire tout seul n'a rien d'irréversible.

\begin{note}[Pourquoi le code reste toujours clair]
Même en thème sombre, la carte du QR reste sur fond blanc. Un code sur fond sombre scanne
moins bien --- la fiabilité du scan prime sur la cohérence du thème, à cet endroit
précisément.
\end{note}

\section{Être prévenu, ajouter en retour}
\label{sec:qr-naif-retour}

Quand quelqu'un utilise votre code --- scan ou lien ---, vous êtes prévenu
immédiatement~: un dialogue propose d'ajouter la personne \textbf{en retour} (« untel
vous a ajouté à ses contacts préférés avec votre code QR. L'ajouter en retour ? »,
\emph{Ajouter} / \emph{Non merci}). Si l'application est en arrière-plan ou fermée, une
notification prend le relais. Et si vous aviez déjà cette personne en contact, pas de
question~: une simple information.

\begin{note}[Le dialogue n'habite aucun écran]
Le code a pu être partagé par lien~: à l'instant où il est utilisé, son propriétaire
peut être n'importe où dans l'application --- pas forcément sur l'écran « Mon code ». Le
dialogue est donc global, affiché où que l'on soit.
\end{note}

L'écran « Mon code », lui, réagit à sa manière~: le code venant d'être consommé, il en
affiche un neuf aussitôt --- le propriétaire peut enchaîner plusieurs personnes sans un
geste.

\section{Reconnaître un contact ajouté par QR}

Un contact ajouté par QR reste reconnaissable après coup, partout où il apparaît~:

\begin{enumerate}[leftmargin=1.6em]
  \item une \textbf{pastille QR} au coin bas-gauche de l'avatar --- le bas-droit reste
        réservé au point de présence en ligne, les deux cohabitent --- dans
        \textbf{toutes} les listes~: contacts préférés, feuille d'ajout, nouvelle
        discussion, sélection de membres, partage de contact\ldots{} et jusqu'à la liste
        des discussions~;
  \item une \textbf{mention datée} sur la fiche du contact~: « Ajouté par QR code ·
        30~juil.~2026 »~;
  \item un \textbf{filtre} « Tous / Par QR » sur l'écran des contacts préférés, chaque
        option affichant son compte.
\end{enumerate}

L'ajout \emph{en retour} est marqué de la même façon~: les deux directions du lien
portent la même origine.

\begin{note}[Maquette interactive]
\href{https://claude.ai/code/artifact/0e095614-eb7c-43bd-8f14-b0d15e61c139}%
{\texttt{claude.ai/code/artifact/0e095614}}\\
Source versionnée~: \code{docs/architecture/maquettes/qr-distinction.html}
\end{note}


La liste des discussions porte le même filtre « Par QR » que l'écran des contacts
préférés~: un chip parmi les filtres existants (Tous, Discussions, Groupes…), qui ne
retient que les conversations avec un contact rencontré par QR.

Enfin, la carte de résultat du scan propose un champ facultatif~: une \textbf{note de
contexte}, quelques mots pour se souvenir de la rencontre. Elle se retrouve sur la fiche
du contact, en italique sous la mention d'origine --- et nulle part ailleurs~: c'est un
aide-mémoire personnel, jamais montré à l'intéressé.

\section{Partager son code}

\textbf{Partager} ouvre un choix explicite~: le \textbf{lien} (cliquable, accompagné de
l'Alanya ID) ou l'\textbf{image} (la carte telle qu'affichée, compte à rebours compris).
Dans les deux cas la légende annonce la couleur --- « Ce code est valable 10~minutes et
pour une seule personne » --- suivie de l'Alanya ID~: le code meurt en dix minutes,
l'Alanya ID reste le chemin permanent, les deux voyagent ensemble.

\begin{note}[Pas d'enregistrement dans la galerie]
Enregistrer la carte en image n'est volontairement pas proposé~: une image morte au bout
de dix minutes finirait en code cassé au fond d'une pellicule. L'option reviendra avec
les codes permanents des comptes business --- eux supportent l'impression et
l'affichage.
\end{note}

\section{Écran --- Scanner}

Plein écran, viseur au centre, façon appareil photo. Dès qu'un code d'identité est
reconnu, une \textbf{carte de résultat} se pose en bas de l'écran~: la personne ajoutée
(avatar, nom), « Ajouté à vos contacts », et trois suites possibles ---
\textbf{Message}, \textbf{Voir détails}, \textbf{Annuler}. La caméra reste vivante
derrière la carte~: on peut enchaîner plusieurs personnes sans ressortir de l'écran. Pas
d'écran intermédiaire, pas de confirmation à taper.

Trois cas particuliers, traités comme des succès doux plutôt que des erreurs~: scanner
son propre code (rien ne se passe, un message discret l'indique --- et le code
\emph{survit}), scanner un contact déjà favori (la carte s'affiche quand même, sans
bouton Annuler~: rien n'a été écrit), scanner un code expiré, déjà utilisé ou illisible
(message clair, la caméra reste active).

Un bouton \textbf{Importer une image} couvre enfin le code reçu en capture d'écran ou en
photo~: l'image choisie dans la galerie est analysée exactement comme le cadre de la
caméra --- même résultat, même carte.

\section{Points d'entrée}

\begin{enumerate}[leftmargin=1.6em]
  \item Le profil, une icône QR à côté de l'Alanya ID
        (\code{lib/screens/profile/profile\_screen.dart}, en-tête).
  \item La feuille d'ajout de contact, un bouton « Scanner un code » au-dessus de la
        recherche existante (\code{lib/widgets/add\_contact\_sheet.dart}).
\end{enumerate}

\chapter{Connexion d'un second appareil}
\label{ch:qr-naif-connexion}

\section{Le principe}

Contrairement à l'ajout de contact, une connexion ne s'approuve \textbf{jamais} au simple
scan. Le nouvel appareil affiche un code ; c'est l'appareil \emph{déjà} connecté qui le
scanne et confirme --- jamais l'inverse.

\begin{alerte}[Pourquoi une confirmation, jamais une reconnaissance instantanée]
Si le scan suffisait à connecter, un attaquant pourrait afficher le QR de connexion d'un
appareil qu'il contrôle --- sur un écran, un site, une affiche --- et attendre qu'une
victime le scanne par réflexe avec son propre compte déjà connecté. C'est une attaque
documentée contre les connexions par QR (le «~quishing~»). L'écran de confirmation, qui
montre \emph{quel} appareil demande la connexion avant de l'accepter, est le seul rempart~:
il doit toujours s'afficher, sans exception ni raccourci.
\end{alerte}

\section{Écran --- Se connecter par QR}

Sur le \textbf{nouvel} appareil, non encore connecté~: un code, un compte à rebours, un
statut (« En attente de scan… »). Le code se régénère seul à l'expiration --- l'utilisateur
n'a jamais un code mort sous les yeux.

\begin{center}
  \imageOuCadre{maquettes/qr-connexion-parcours.png}{\linewidth}%
    {La poignée de main~: le nouvel appareil affiche, l'appareil de confiance confirme\\
     \texttt{maquettes/qr-connexion-parcours.png}}
\end{center}

\begin{note}[Maquette interactive]
\href{https://claude.ai/code/artifact/a69b2050-b2dc-4aa1-a204-59ada2526a86}%
{\texttt{claude.ai/code/artifact/a69b2050}}\\
Source versionnée~: \code{docs/architecture/maquettes/qr-connexion.html}
\end{note}

\section{Écran --- Confirmer la connexion}

Sur l'appareil \textbf{déjà} connecté, après le scan~: une fiche décrivant ce qui vient
d'être scanné --- type d'appareil, ville approximative, à l'instant --- suivie de deux
boutons, \textbf{Refuser} et \textbf{Confirmer}. Un avertissement rappelle de refuser et de
changer son mot de passe en cas de doute.

\begin{note}[Pourquoi l'adresse compte plus que le nom de l'appareil]
Le nom et la plateforme sont \emph{annoncés} par l'appareil qui demande la connexion~: rien
n'empêche un attaquant d'appeler le sien « Alanya --- Vérification de sécurité ». La seule
information que le serveur \emph{constate} lui-même est l'adresse d'où part la demande.
C'est donc la seule sur laquelle fonder sa décision, et l'écran le dit explicitement.
\end{note}

\begin{note}[Une ville, pas une adresse technique]
Cette adresse est traduite en lieu approximatif --- « Douala, Cameroun » --- avant d'être
montrée. Sous sa forme brute (\code{41.202.219.14}), elle n'apprendrait rien à qui n'est pas
technicien~: on afficherait la seule information fiable de l'écran sous la seule forme où
elle est inutilisable. Sous forme de ville, le doute devient une évidence --- on sait
immédiatement si la demande part d'où l'on se trouve. Si la traduction n'aboutit pas,
l'adresse brute reste affichée plutôt que rien.
\end{note}

\section{Écran --- Mes appareils connectés}

Un écran à part, indépendant de toute connexion en cours, atteint depuis \emph{Compte et
sécurité}~: tous les appareils sur lesquels le compte est ouvert --- pas seulement ceux
ajoutés par QR --- avec un repère « Cet appareil » et un bouton \textbf{Déconnecter} par
ligne, à effet immédiat.

\begin{center}
  \imageOuCadre{maquettes/qr-connexion-appareils.png}{\linewidth}%
    {Mes appareils connectés~: tous les logins, pas seulement ceux par QR\\
     \texttt{maquettes/qr-connexion-appareils.png}}
\end{center}

\begin{note}[Maquette interactive]
\href{https://claude.ai/code/artifact/a69b2050-b2dc-4aa1-a204-59ada2526a86}%
{\texttt{claude.ai/code/artifact/a69b2050}} --- second panneau de la même maquette\\
Source versionnée~: \code{docs/architecture/maquettes/qr-connexion.html}
\end{note}

\section{Points d'entrée}

\begin{enumerate}[leftmargin=1.6em]
  \item L'écran de connexion, un bouton « Se connecter avec un code QR »
        (\code{screens/authentification/\allowbreak login\_screen.dart}), pour l'appareil
        qui n'est pas encore connecté.
  \item \emph{Compte et sécurité}, une nouvelle section « Appareils connectés »
        (\code{screens/profile/\allowbreak account\_security\_\allowbreak screen.dart}),
        pour lier un nouvel appareil ou consulter la liste depuis un appareil déjà
        connecté.
\end{enumerate}
